\section{Verification-Aware Decision Rule}\label{sec:mp_decision}

After each paid target batch, the posterior supplies a predictive mean and
covariance for $(f,g,\eta)$ over the public candidate bank.  A candidate's
deployment value is the expected objective quality multiplied by the
probability that an independent verifier will certify its constraint.  For a
fixed recommendation $x$, the verifier uses fresh replications and the frozen
chance level $\alpha$; search observations never enter this test.

The acquisition rule evaluates the expected increase in this verified utility
under the remaining target and verification budget.  It uses the correlated
Gaussian approximation for the joint mean--risk posterior and finite
quadrature for the fantasy observation.  The live physical-risk posterior updates factor-family weights; the mean
function posteriors are updated by Gaussian conditioning.  The fantasy
calculation uses the frozen approximation analyzed in the
A45 numerical diagnostic.  That diagnostic found a maximum 9-versus-13-node
difference of $2.93\times10^{-5}$ in the controlled risk-only probe, so the
existing acquisition rule is retained for the reported pilots.

This utility is a decision rule, not a theorem that the selected policy is the
global optimum of the chance-constrained problem.  We report the initial
recommendation, final recommendation, true post-decision reward, independent
certificate, and false-certificate indicator separately.  A successful
certificate requires a fresh verifier; a posterior feasibility probability is
never reported as a deployment certificate.
